% GENERATED FILE. Edit canonical lessons, then run npm run generate:books.
% canonical-source-hash: fnv1a64:02eee47d775a70b4
% canonical-lessons: TE-C349-nyayasthanam, TE-C349-nyayamurti, TE-C349-nyayavadi, TE-C349-tirpu, TE-C349-beyilu, TE-R349-review-342-349

\chapter{In Court}
\label{ch:te-in-court}
\hlchaptermodality{349}

\begin{chapteropening}
\textbf{By the end of this chapter:} \emph{I can say a court of law, a judge, an advocate, a verdict and bail.}

Complete the last lesson of chapter 349: I can say a court of law, a judge, an advocate, a verdict and bail.
\end{chapteropening}

\section[\texorpdfstring{\te{న్యాయస్థానం} (nyāyasthānaṁ)}{న్యాయస్థానం (nyāyasthānaṁ)}]{\te{న్యాయస్థానం} (nyāyasthānaṁ) — a court of law}

\label{lesson:TE-C349-nyayasthanam}



\begin{quote}
Before the new one: say the Telugu for to lie in wait, then the Telugu for a disguise.
\end{quote}

\subsection*{You'll want to know: \te{న్యాయస్థానం}}
\textbf{\te{న్యాయస్థానం}} — \emph{nyāyasthānaṁ} — \textquotedblleft{}a court of law\textquotedblright{}.

\begin{quote}
\te{ఆ} \te{కేసు} \te{న్యాయస్థానంలో} \te{ఉంది}.
\end{quote}

\begin{quote}
\emph{ā kēsu nyāyasthānaṁlō undi.}
\end{quote}

\begin{quote}
That case is in court.
\end{quote}

\te{న్యాయం}, justice, and \te{స్థానం}, a place; in speech \te{కోర్టు}.

\begin{cousinweb}
\textbf{\te{న్యాయస్థానం}} (\emph{nyāyasthānaṁ}) is a compound borrowed from Sanskrit: \emph{nyāya}, \textquotedblleft{}rule, method; justice\textquotedblright{} (CDIAL 7614), and \emph{sthāna}, \textquotedblleft{}a place\textquotedblright{} (CDIAL 13753).
\end{cousinweb}

\begin{grammarlens}[title={What you've built}]
One more word for telling what happened.
\end{grammarlens}

\subsection*{Your turn}
\begin{itemize}
\raggedright
  \item \emph{Say it:} \emph{nyāyasthānaṁ}
  \item \emph{Say it:} \emph{nyāyasthānaṁ}, once more
  \item \emph{Say it:} \emph{ā kēsu nyāyasthānaṁlō undi.} — That case is in court.
  \item \emph{Recall:} say the Telugu for to lie in wait, then the Telugu for a disguise, then say \emph{nyāyasthānaṁ} again
\end{itemize}

\begin{tcolorbox}[breakable,colback=teal!4,colframe=teal!35!black,title={Before you move on}]
What does \te{న్యాయస్థానం} mean? (\textquotedblleft{}A court of law\textquotedblright{}.) Say it once more.
\end{tcolorbox}
\section[\texorpdfstring{\te{న్యాయమూర్తి} (nyāyamūrti)}{న్యాయమూర్తి (nyāyamūrti)}]{\te{న్యాయమూర్తి} (nyāyamūrti) — a judge}

\label{lesson:TE-C349-nyayamurti}



\begin{quote}
Before the new one: say the Telugu for a disguise, then the Telugu for a court of law.
\end{quote}

\subsection*{You'll want to know: \te{న్యాయమూర్తి}}
\textbf{\te{న్యాయమూర్తి}} — \emph{nyāyamūrti} — \textquotedblleft{}a judge\textquotedblright{}.

\begin{quote}
\te{న్యాయమూర్తి} \te{సాక్షిని} \te{చాలా} \te{ప్రశ్నలు} \te{అడిగాడు}.
\end{quote}

\begin{quote}
\emph{nyāyamūrti sākṣini cālā praśnalu aḍigāḍu.}
\end{quote}

\begin{quote}
The judge asked the witness many questions.
\end{quote}

Plural: \te{న్యాయమూర్తులు}, \emph{nyāyamūrtulu}.

\te{న్యాయం}, justice, and \te{మూర్తి}, a form: justice in person.

\begin{cousinweb}
\textbf{\te{న్యాయమూర్తి}} (\emph{nyāyamūrti}) is borrowed from Sanskrit: \emph{nyāya}, \textquotedblleft{}rule; justice\textquotedblright{} (CDIAL 7614), and \emph{mūrti}, \textquotedblleft{}a figure, a form\textquotedblright{} (CDIAL 10245).
\end{cousinweb}

\begin{grammarlens}[title={What you've built}]
One more word for telling what happened.
\end{grammarlens}

\subsection*{Your turn}
\begin{itemize}
\raggedright
  \item \emph{Say it:} \emph{nyāyamūrti}
  \item \emph{Say it:} \emph{nyāyamūrti}, once more
  \item \emph{Say it:} \emph{nyāyamūrti sākṣini cālā praśnalu aḍigāḍu.} — The judge asked the witness many questions.
  \item \emph{Recall:} say the Telugu for a disguise, then the Telugu for a court of law, then say \emph{nyāyamūrti} again
\end{itemize}

\begin{tcolorbox}[breakable,colback=teal!4,colframe=teal!35!black,title={Before you move on}]
What does \te{న్యాయమూర్తి} mean? (\textquotedblleft{}A judge\textquotedblright{}.) Say it once more.
\end{tcolorbox}
\section[\texorpdfstring{\te{న్యాయవాది} (nyāyavādi)}{న్యాయవాది (nyāyavādi)}]{\te{న్యాయవాది} (nyāyavādi) — an advocate; a lawyer}

\label{lesson:TE-C349-nyayavadi}



\begin{quote}
Before the new one: say the Telugu for a court of law, then the Telugu for a judge.
\end{quote}

\subsection*{You'll want to know: \te{న్యాయవాది}}
\textbf{\te{న్యాయవాది}} — \emph{nyāyavādi} — \textquotedblleft{}an advocate; a lawyer\textquotedblright{}.

\begin{quote}
\te{న్యాయవాది} \te{నిందితుడి} \te{తరఫున} \te{మాట్లాడాడు}.
\end{quote}

\begin{quote}
\emph{nyāyavādi nindituḍi taraphuna māṭlāḍāḍu.}
\end{quote}

\begin{quote}
The advocate spoke on behalf of the accused.
\end{quote}

Plural: \te{న్యాయవాదులు}, \emph{nyāyavādulu}.

A formal word for \te{వకీలు}: one who argues for justice.

\begin{cousinweb}
\textbf{\te{న్యాయవాది}} (\emph{nyāyavādi}) is borrowed from Sanskrit \emph{nyāyavādin}, \textquotedblleft{}speaking properly\textquotedblright{} (Monier-Williams), from \emph{nyāya}, \textquotedblleft{}rule, justice\textquotedblright{}, and \emph{vādin}, \textquotedblleft{}speaking\textquotedblright{}.
\end{cousinweb}

\begin{grammarlens}[title={What you've built}]
One more word for telling what happened.
\end{grammarlens}

\subsection*{Your turn}
\begin{itemize}
\raggedright
  \item \emph{Say it:} \emph{nyāyavādi}
  \item \emph{Say it:} \emph{nyāyavādi}, once more
  \item \emph{Say it:} \emph{nyāyavādi nindituḍi taraphuna māṭlāḍāḍu.} — The advocate spoke on behalf of the accused.
  \item \emph{Recall:} say the Telugu for a court of law, then the Telugu for a judge, then say \emph{nyāyavādi} again
\end{itemize}

\begin{tcolorbox}[breakable,colback=teal!4,colframe=teal!35!black,title={Before you move on}]
What does \te{న్యాయవాది} mean? (\textquotedblleft{}An advocate; a lawyer\textquotedblright{}.) Say it once more.
\end{tcolorbox}
\section[\texorpdfstring{\te{తీర్పు} (tīrpu)}{తీర్పు (tīrpu)}]{\te{తీర్పు} (tīrpu) — a verdict; a judgment}

\label{lesson:TE-C349-tirpu}



\begin{quote}
Before the new one: say the Telugu for a judge, then the Telugu for an advocate.
\end{quote}

\subsection*{You'll want to know: \te{తీర్పు}}
\textbf{\te{తీర్పు}} — \emph{tīrpu} — \textquotedblleft{}a verdict; a judgment\textquotedblright{}.

\begin{quote}
\te{న్యాయమూర్తి} \te{నిన్న} \te{తీర్పు} \te{చెప్పాడు}.
\end{quote}

\begin{quote}
\emph{nyāyamūrti ninna tīrpu ceppāḍu.}
\end{quote}

\begin{quote}
The judge gave the verdict yesterday.
\end{quote}

\te{తీర్పు} \te{చెప్పు}, to pronounce a verdict; it comes from \te{తీరు}, to be settled.

\begin{cousinweb}
\textbf{\te{తీర్పు}} (\emph{tīrpu}) is native Dravidian: DEDR 3278 lists Telugu \emph{tīrpu}, a settling, beside \emph{tīru}, \textquotedblleft{}to be finished, set right\textquotedblright{}.
\end{cousinweb}

\begin{grammarlens}[title={What you've built}]
One more word for telling what happened.
\end{grammarlens}

\subsection*{Your turn}
\begin{itemize}
\raggedright
  \item \emph{Say it:} \emph{tīrpu}
  \item \emph{Say it:} \emph{tīrpu}, once more
  \item \emph{Say it:} \emph{nyāyamūrti ninna tīrpu ceppāḍu.} — The judge gave the verdict yesterday.
  \item \emph{Recall:} say the Telugu for a judge, then the Telugu for an advocate, then say \emph{tīrpu} again
\end{itemize}

\begin{tcolorbox}[breakable,colback=teal!4,colframe=teal!35!black,title={Before you move on}]
What does \te{తీర్పు} mean? (\textquotedblleft{}A verdict; a judgment\textquotedblright{}.) Say it once more.
\end{tcolorbox}
\section[\texorpdfstring{\te{బెయిలు} (beyilu)}{బెయిలు (beyilu)}]{\te{బెయిలు} (beyilu) — bail}

\label{lesson:TE-C349-beyilu}



\begin{quote}
Before the new one: say the Telugu for an advocate, then the Telugu for a verdict.
\end{quote}

\subsection*{You'll want to know: \te{బెయిలు}}
\textbf{\te{బెయిలు}} — \emph{beyilu} — \textquotedblleft{}bail\textquotedblright{}.

\begin{quote}
\te{నిందితుడికి} \te{బెయిలు} \te{దొరికింది}.
\end{quote}

\begin{quote}
\emph{nindituḍiki beyilu dorikindi.}
\end{quote}

\begin{quote}
The accused got bail.
\end{quote}

\te{బెయిలు} \te{మీద} \te{విడుదల}, release on bail.

\begin{cousinweb}
\textbf{\te{బెయిలు}} (\emph{beyilu}) is a loan from English \emph{bail}, with the final -\te{ఉ}.
\end{cousinweb}

\begin{grammarlens}[title={What you've built}]
One more word for telling what happened.
\end{grammarlens}

\subsection*{Your turn}
\begin{itemize}
\raggedright
  \item \emph{Say it:} \emph{beyilu}
  \item \emph{Say it:} \emph{beyilu}, once more
  \item \emph{Say it:} \emph{nindituḍiki beyilu dorikindi.} — The accused got bail.
  \item \emph{Recall:} say the Telugu for an advocate, then the Telugu for a verdict, then say \emph{beyilu} again
\end{itemize}

\begin{tcolorbox}[breakable,colback=teal!4,colframe=teal!35!black,title={Before you move on}]
What does \te{బెయిలు} mean? (\textquotedblleft{}Bail\textquotedblright{}.) Say it once more.
\end{tcolorbox}
\section[(dialogue)]{Words from chapters 342–349 — a review}

\label{lesson:TE-R349-review-342-349}



\begin{quote}
Say the Telugu for a drum proclamation and a well-known person.
\end{quote}

\subsection*{Your turn}
Say these aloud:

\begin{itemize}
\raggedright
  \item a drum proclamation, a well-known person, a public meeting, a speech, to give a speech
  \item the injured, a tragic accident, a stampede, narrowly, to get out
  \item condolence, death, one's last breath, a funeral procession, a funeral
  \item a cremation ground, a grave, a tribute, a death anniversary, a will
  \item an engagement, a marriage, a wedding invitation, a wedding pendant, to invite
  \item a feast, a present, a bouquet, to greet warmly, to meet
  \item a robbery, to attack, to snatch, to lie in wait, a disguise
  \item a court of law, a judge, an advocate, a verdict, bail
\end{itemize}

\begin{tcolorbox}[breakable,colback=teal!4,colframe=teal!35!black,title={Before you move on}]
A drum proclamation? (\textbf{\te{దండోరా}}, \emph{daṇḍōrā}.) A well-known person? (\textbf{\te{ప్రముఖుడు}}, \emph{pramukhuḍu}.) A public meeting? (\textbf{\te{సభ}}, \emph{sabha}.) A speech? (\textbf{\te{ప్రసంగం}}, \emph{prasaṅgaṁ}.) To give a speech? (\textbf{\te{ప్రసంగించు}}, \emph{prasaṅgiñcu}.) The injured? (\textbf{\te{క్షతగాత్రులు}}, \emph{kṣatagātrulu}.) A tragic accident? (\textbf{\te{దుర్ఘటన}}, \emph{durghaṭana}.) A stampede? (\textbf{\te{తొక్కిసలాట}}, \emph{tokkisalāṭa}.) Narrowly? (\textbf{\te{త్రుటిలో}}, \emph{truṭilō}.) To get out? (\textbf{\te{బయటపడు}}, \emph{bayaṭapaḍu}.) Condolence? (\textbf{\te{సంతాపం}}, \emph{santāpaṁ}.) Death? (\textbf{\te{మరణం}}, \emph{maraṇaṁ}.) One's last breath? (\textbf{\te{తుదిశ్వాస}}, \emph{tudiśvāsa}.) A funeral procession? (\textbf{\te{అంతిమయాత్ర}}, \emph{antimayātra}.) A funeral? (\textbf{\te{అంత్యక్రియలు}}, \emph{antyakriyalu}.) A cremation ground? (\textbf{\te{శ్మశానం}}, \emph{śmaśānaṁ}.) A grave? (\textbf{\te{సమాధి}}, \emph{samādhi}.) A tribute? (\textbf{\te{నివాళి}}, \emph{nivāḷi}.) A death anniversary? (\textbf{\te{వర్ధంతి}}, \emph{vardhanti}.) A will? (\textbf{\te{వీలునామా}}, \emph{vīlunāmā}.) An engagement? (\textbf{\te{నిశ్చితార్థం}}, \emph{niścitārthaṁ}.) A marriage? (\textbf{\te{వివాహం}}, \emph{vivāhaṁ}.) A wedding invitation? (\textbf{\te{శుభలేఖ}}, \emph{śubhalēkha}.) A wedding pendant? (\textbf{\te{తాళి}}, \emph{tāḷi}.) To invite? (\textbf{\te{ఆహ్వానించు}}, \emph{āhvāniñcu}.) A feast? (\textbf{\te{విందు}}, \emph{vindu}.) A present? (\textbf{\te{కానుక}}, \emph{kānuka}.) A bouquet? (\textbf{\te{పుష్పగుచ్ఛం}}, \emph{puṣpagucchaṁ}.) To greet warmly? (\textbf{\te{స్వాగతించు}}, \emph{svāgatiñcu}.) To meet? (\textbf{\te{కలుసుకో}}, \emph{kalusukō}.) A robbery? (\textbf{\te{దోపిడీ}}, \emph{dōpiḍī}.) To attack? (\textbf{\te{దాడి} \te{చేయు}}, \emph{dāḍi cēyu}.) To snatch? (\textbf{\te{లాక్కో}}, \emph{lākkō}.) To lie in wait? (\textbf{\te{కాపుకాచు}}, \emph{kāpukācu}.) A disguise? (\textbf{\te{మారువేషం}}, \emph{māruvēṣaṁ}.) A court of law? (\textbf{\te{న్యాయస్థానం}}, \emph{nyāyasthānaṁ}.) A judge? (\textbf{\te{న్యాయమూర్తి}}, \emph{nyāyamūrti}.) An advocate? (\textbf{\te{న్యాయవాది}}, \emph{nyāyavādi}.) A verdict? (\textbf{\te{తీర్పు}}, \emph{tīrpu}.) Bail? (\textbf{\te{బెయిలు}}, \emph{beyilu}.)
\end{tcolorbox}
